\section{Discussion}

We observe consistent, significant gains  with self-critique strategies over baselines.
\Cref{tab:aggregate_gain} shows that \EPSEMCRIT{} improves over \EPLABEL{} on 32 of 42 (model, dataset) points. The remaining 10 non-improving cases are not randomly distributed; they cluster into three interpretable groups, each with an independent explanation.

\textbf{Llama 4 Scout} Four non-improving cases involve Llama 4 Scout (Multi-Condition Ranking, PubMed, Steam Pref, and Anime Pref). This is consistent with its low average suggestibility of 56.2\% (\Cref{tab:suggestibility}), which limits its ability to benefit from any in-context guidance, critiques included. Low suggestibility appears to impose a ceiling on gains from memory augmentation regardless of critique quality.

\textbf{Reasoning models on preference data} Four cases involve reasoning models on preference datasets: GPT-5 on Anime Pref and Movie Pref, and Qwen3-thinking on Anime Pref and Book Pref (the latter a tie). On these datasets \EPLABEL{} already yields substantial gains over \ZEROSHOT{} --- 8-20 absolute points across the affected model-dataset pairs --- and our token analysis in \Cref{sec:token_analysis} shows both models reduce thinking tokens substantially when critiques are provided. Taken together, this suggests these reasoning models perform critique-like reasoning internally over retrieved examples, making explicit critiques largely redundant for accuracy while still delivering the efficiency benefits documented in \Cref{sec:token_analysis}.

\textbf{PubMed} The remaining two cases are both on PubMed, for Qwen3-instruct and Qwen3-thinking, consistent with the pattern identified in \Cref{sec:main_results}: PubMed saw the least improvement across all models compared with both \ZEROSHOT{} and \EPLABEL{}. PubMed requires very specific factual knowledge that is rarely shared across questions, limiting how much critiques generated on training examples can transfer to test questions.

Each of the three groups points to a different boundary condition on the method: model suggestibility limitations, high token limits that permit internal critique-like reasoning, and domain structure that limits cross-instance transfer of critiques. None of the failures reflect random noise.

A second major pattern we observed was the consistent preference for episodic over semantic memory. This reflects a fundamental distinction between lazy and eager generalization strategies. Much like k-nearest neighbor methods, which often outperform regression-based approaches when training data is abundant, our lazy-learning episodic method has the advantage of not needing to commit to learning the full function over the entire domain. Instead, it learns only a local approximation to the function at the current query point, allowing for greater flexibility and specificity that may be lost during semantic abstraction.

Our analysis has centered on the accuracy of different agentic learning strategies, but design choices also impact computational cost and the ability to incorporate ongoing supervision. Semantic memory requires additional cost and complexity over episodic memory. At inference time, however, semantic memory offers more readily applicable knowledge, while episodic memory relies on retrieval quality. This trade-off suggests that the optimal strategy may depend on the size of the supervised dataset and the frequency of inference---semantic memory may be better suited for frequent inference under sparse supervision, while episodic memory may be preferable when supervision is abundant and retrieval is reliable.

An additional interesting direction for exploring model suggestibility is to disentangle how much a model's behavior changes due to genuinely incorporating supervised signals into its internal beliefs versus merely adapting its responses to please the user. In our empirical study, we observed that models exhibited higher suggestibility scores when critiques were attributed to the user, compared to when the same critiques were believed to originate from the model itself or another model. This suggests that the perceived source of feedback plays a significant role in how seriously the model treats the signal, opening up opportunities to better understand and guide belief formation in interactive learning systems.

\section{Conclusion}

In this work, we introduced a reflective framework for agent adaptation that leverages episodic and semantic memory to learn from labeled supervision without parameter updates. Our experiments demonstrate that by storing and retrieving label-driven critiques, LLMs can achieve significant performance gains. Our best-performing strategy, \EPSEMCRIT, yields an average improvement of 8.1 percentage points over zero-shot baselines and 4.6 points over RAG-style few-shot baselines.

Beyond accuracy, our analysis through the lens of suggestibility reveals that a model's capacity to utilize self-reflection is a function of its receptivity to external reasoning, rather than just its scale. Furthermore, we demonstrate a critical efficiency benefit: reasoning models can 'offload' internal computation to retrieved critiques, reducing thinking tokens by an average of 31.95\%.

This memory-driven approach offers a lightweight, interpretable, and computationally efficient alternative to fine-tuning. By treating model adaptation as a reflective process rather than a purely statistical one, we move closer to agents that can learn from their experiences in real-time, maintaining high performance across diverse and evolving tasks.